\documentclass{article}
\usepackage[preprint]{corl_2026} % submission: \usepackage{corl_2026}, which anonymises the authors; final: [final]
\usepackage{graphicx}
\usepackage{booktabs}
\usepackage{amsmath}

\newif\ifdraftmarks \draftmarkstrue
\newif\ifwithappendix \withappendixfalse % submission: body plus references; \withappendixtrue for the camera-ready or an anonymous supplement
\newcommand{\appref}[1]{\ifwithappendix Appendix~\ref{#1}\else the supplementary material\fi}
\newcommand{\tbd}{\ifdraftmarks\textcolor[rgb]{0,0.3,0.75}{[tbd]}\else[tbd]\fi}
\newcommand{\todo}[1]{\ifdraftmarks\textcolor[rgb]{0,0.3,0.75}{[#1]}\else[#1]\fi}
\newcommand{\prov}[1]{\ifdraftmarks\textcolor[rgb]{0,0.3,0.75}{#1}\else#1\fi}
% \figslot{file}{width}{height}{what it will show}: the figure when its file exists, a labelled box otherwise
\newcommand{\figslot}[4]{\IfFileExists{#1}{\includegraphics[width=#2,height=#3,keepaspectratio]{#1}}{\fbox{\parbox[c][\dimexpr#3-2\fboxsep-2\fboxrule\relax][c]{\dimexpr#2-2\fboxsep-2\fboxrule\relax}{\centering\footnotesize\textbf{#4}\\[2pt]{\tiny\nolinkurl{#1}}\\[1pt]\textcolor[rgb]{0.8,0,0}{pending}}}}}

\title{DoomShift: Efficient Adaptation of World Models to Domain Shifts}

\author{
  Rohan Nagabhirava\\
  Carnegie Mellon University\\
  \texttt{rnagabhi@andrew.cmu.edu}
  \And
  Keerthana Chirumamilla\\
  Carnegie Mellon University\\
  \texttt{kchiruma@andrew.cmu.edu}
  \And
  Changliu Liu\\
  Carnegie Mellon University\\
  \texttt{cliu6@andrew.cmu.edu}
}

\begin{document}
\maketitle

%===============================================================================

\begin{abstract}
World models have become capable simulators, but they are tested primarily in the domain they were trained on.
A small domain shift, a new arena of the same game, makes their outputs drift and lose realism, though the physics is unchanged and only the situation differs.
DoomShift asks how quickly and cheaply a world model can be adapted across such a shift: one Doom level pack has 17 arenas; we train three pretrained diffusion backbones on four and score each of the other thirteen against persistence (the last frame copied forward), zero-shot for all three and after adaptation for the U-Net.
Off the training maps the U-Net still beats persistence on every arena, but its advantage shrinks from 5.1~dB in distribution to 1.1--2.9~dB (median 2.0) and its LPIPS is worse than persistence's on 13 of 13, while its turn response survives.
Rank-16 adapters trained on eight episodes close half of an arena's gap to the in-distribution advantage within 4k updates on 9 of 13 arenas, recover a median 1.8~dB and beat persistence perceptually on 11 of 13.
\end{abstract}

\keywords{World models, Domain shift, Few-shot adaptation}

%===============================================================================

\section{Introduction}
\label{sec:intro}

Game world models now render playable worlds from pixels and controls.
GameNGen runs Doom inside a diffusion model~\citep{valevski2024gamengen}, DIAMOND learns Atari and Counter-Strike in pixel space~\citep{alonso2024diamond} and Matrix-Game~2 scales the recipe to open worlds~\citep{he2025matrixgame2}.
Each of these models is evaluated only in the environment it was trained on: the reported score is a held-out trajectory drawn from its own training scenes and none has been shown to generalize to a new environment within the same game.
This holds even for MultiGen, which trains a single Doom model across 100 procedurally generated maps with a pre-trained agent~\citep{po2026multigen}: it reports PSNR, SSIM and LPIPS only on its training distribution, with no map held out for evaluation.

A learned simulator is useful only where it is faithful and a robot or an agent moving to a new room, level or workspace needs to know how much of its world model still holds there and what it costs to make the rest hold.
Game worlds let us pose this question with the engine held fixed and only the environment changed: a new arena of the same game preserves the physics and the controls exactly and only the layout, textures and lighting differ.
Yet unchanged dynamics alone do not guarantee a faithful render: pretrained video and world models are typically moved to a new domain only through a further full fine-tune~\citep{gao2025adaworld} or a dedicated adapter on the frozen backbone~\citep{rigter2024avid} and a model trained on a handful of maps tends to fall back on the appearance it already knows once the surroundings change.
Training a separate model per map is prohibitively costly, so the open question is how cheaply post-training can close this gap once it appears: when a world model meets a new arena of the same game, what does it retain, what does it lose and how many episodes and updates does it take to relearn the rest.

Our paper addresses these concerns and contributes the following:
\begin{itemize}\setlength{\itemsep}{0pt}
\item \textbf{DoomShift}, a benchmark that turns this question into a concrete testbed: 17 arenas of one Doom level pack, three pretrained diffusion backbones trained on four of them and the remaining thirteen scored both zero-shot and after post-training. We release every-tic recordings of \prov{6,000} episodes on the four training maps and 24 on each of the 13 unseen arenas (this paper trains on 2,000 of them), along with the three trained models, the decoder fine-tune and the adapter path (Figure~\ref{fig:method}; link withheld for review).
\item A post-training recipe that adapts a frozen model to a new arena from just eight episodes with a small adapter, together with the per-arena protocol that scores the budget it requires.
\item Evidence that this generalizes without heavy retraining: half of the gap to the training maps is closed on 9 of 13 unseen arenas within 4k updates, most of it within 250 and a perceptual advantage over persistence (the last frame copied forward) is restored on 11 of 13 (Figure~\ref{fig:adapt}; full fine-tune comparator \tbd{}).
\item A characterization of the shift and what predicts it: on unseen arenas the PSNR advantage over persistence shrinks and the perceptual margin changes sign for all three backbones, while the turn response survives (Figures~\ref{fig:teaser} and~\ref{fig:shift}). The model's own zero-shot skill score turns out to be the strongest predictor of where an arena ends up, while the registered distances between recordings order nothing.
\end{itemize}

\begin{figure}[t]
\centering
\figslot{figures/fig_teaser_B.pdf}{\linewidth}{1.5in}{Figure 1: teaser, two rows}
\caption{On an unseen arena the U-Net keeps moving and firing but renders the arena in training-map masonry; eight episodes restore it. Left, training maps 5 and 3 (in distribution); right, unseen arena 7 zero-shot and after 4k adapter updates on 8 episodes. Selected windows, rolled out 16 tics past the context frame with the row's button held throughout; fine-tuned decoder, raw ground truth, scene PSNR; persistence scores 16 to 19~dB on these windows.}
\label{fig:teaser}
\end{figure}

%===============================================================================

\section{Prior work}
\label{sec:related}

Most Doom and game world models are trained on a set of maps and tested on maps from that same set. Papers that test on maps the model never trained on exist, but they keep aside something other than a map from the same game: a new physical body for the agent (concurrent work~\citep{chen2026xeworld}), a whole new game or environment~\citep{rigter2024avid,gao2025adaworld}, or a new building~\citep{koh2021pathdreamer}. Doom itself is tested this way, but only for the agents that play it~\citep{lample2017arnold,wydmuch2018vizdoom}, not for the world models built to simulate it. No prior paper tests a Doom world model on a map from its own game that it never trained on.

We take our method from two papers that move a pretrained video model into a new setting. AVID shows that training a small adapter on top of a frozen model is enough to do this~\citep{rigter2024avid}; Vista uses the same frozen-model-plus-adapter idea to add action control rather than move the model to new content~\citep{gao2024vista}. Neither tests this on a map kept aside from the same game and neither reports what adapting actually costs: how many episodes, how many training updates, for one new scene. AdaWorld solves a similar problem with a short full fine-tune instead of a small adapter~\citep{gao2025adaworld}; we use it as the full-fine-tune comparison our own method should eventually be measured against. DoomShift combines what none of these do: maps kept aside from within a single game, with the adaptation cost measured directly. As our comparison for held-out maps we use persistence, which is using the last frame forward, the standard reference in video prediction~\citep{mathieu2016deep,lotter2017prednet,villegas2019fidelity} that game world models have dropped in favor of raw PSNR, a number that mostly measures how still a video is rather than how good the prediction is. We also compare every score to the reconstruction upper bound: encoding and decoding the true next frame and measuring the quality lost in that round trip alone~\citep{zheng2023occworld,karypidis2024dinoforesight}, which tells us how much error is the model's fault versus a limit built into its own decoder.
%===============================================================================

\section{Method}
\label{sec:bench}

\textbf{Data.} We collect the data for this research by running experiments with the Arnold agent~\citep{lample2017arnold} through ViZDoom~\citep{kempka2016vizdoom} deathmatches: free-for-all matches where each player scores by killing the others. Arnold's level pack has 17 arenas. A level pack is a single Doom WAD file that bundles a set of maps and each arena is one of those maps, with its own layout, textures and lighting. For every episode, we let Arnold play against 8 bots for 150 seconds. We record every tic, the engine's 1/35~s simulation step, capturing a 320$\times$240 frame and the 19 buttons Arnold pressed. We split the 17 arenas into two groups. Four arenas (2 to 5) are our training maps: we record 500 episodes on each, then hold out 25 validation episodes per map (512 windows) for scoring. The other thirteen arenas (1, 6 to 17) are never trained on: we record 24 new episodes on each, using 16 for adaptation and holding out 8 (256 windows) for scoring. In total we collect 2,412 episodes, about 96 hours of play and about 12.1 million frames.

\textbf{Three backbones, one training recipe.} We pick three pretrained backbones to ask three different questions. Our anchor is the SD~1.4 U-Net~\citep{rombach2022ldm} (860M parameters): it is GameNGen's backbone~\citep{valevski2024gamengen}, so our in-distribution numbers come from the same network as the published Doom result. To test how different architectures perform at fixed scale, we swap in PixArt-$\alpha$~\citep{chen2023pixart}, a 628M transformer in the same latent space instead of a U-Net; architecture turns out not to change what survives the shift (in distribution their PSNR and $A$ differ by under 0.1~dB; on unseen arenas their median $A$ falls to 2.03 and 2.19~dB and both lose to persistence in LPIPS on all 13; Table~\ref{tab:unseen}). To test the performace with scaling of the model, we use SD~3.5 Medium~\citep{esser2024sd3}, a widely adopted open MMDiT (2.27B parameters, 16-channel latents); its latent space changes along with its scale, so this step compares pretrained packages rather than architectures.

We turn all three into next-tic world models under one shared recipe (Figure~\ref{fig:method}). We give each model 32 context tics so it can see motion, train with $v$-prediction~\citep{salimans2022progressive} for stable training across noise levels and add context noise augmentation as in GameNGen~\citep{valevski2024gamengen} so the model tolerates its own rollout errors. We train each backbone for 200k updates at batch size 32, a step count we match across backbones, keep an EMA at 0.9999 in fp32 (a bf16 EMA underflows and freezes) and sample with 10-step DDIM~\citep{song2021ddim} to keep evaluation cheap (\appref{app:protocol}).

\textbf{Metrics} We score every predicted frame $\hat z$ against four questions: is it more accurate than just copying the last frame forward, does it look more realistic than that copy, how far is it from the best the decoder could possibly render and does it turn the right way when we mirror the controls? We compute each score using the ground-truth next frame $x$ (latent $z$), the last context frame $x_\text{last}$ (latent $z_\text{last}$) and the decoder $D$.

For accuracy, we compute the \emph{decoded advantage} $A=\mathrm{PSNR}(D(\hat z),D(z))-\mathrm{PSNR}(D(z_\text{last}),D(z))$, a paired difference in the spirit of Genie's $\Delta$PSNR~\citep{bruce2024genie}, but using persistence rather than random actions as the reference. For realism, we compute the \emph{perceptual margin} $M=\mathrm{LPIPS}(D(\hat z),x)-\mathrm{LPIPS}(x_\text{last},x)$~\citep{zhang2018lpips}, which is negative exactly when the prediction looks better than the copy. For distance from the decoder's ceiling, we compute the \emph{gap} $G=\mathrm{PSNR}(D(z),x)-\mathrm{PSNR}(D(\hat z),x)$ to the reconstruction upper bound. For turning, we compute the \emph{directional} score: the fraction of turning windows whose predicted view turns the other way once we swap left and right in the newest control.

We also ask a fifth question before any decoding happens: how good is the model in its own latent space? We answer it with the \emph{zero-shot latent skill score} $S_0$, $10\log_{10}$ of persistence's latent MSE over the model's. We score $M$ and $G$ against the raw frame and restrict pixel scores to the scene rows 0 to 207, since persistence copies the HUD almost exactly. We use $A_\text{train}$, the advantage on the training maps (5.06~dB, episode-bootstrap interval 4.75 to 5.38), as the in-distribution reference every arena is measured against, not as an arena's own upper bound.

\textbf{Training setup.} We train each backbone on a single RTX A6000 (49~GB) in bf16 with fused AdamW, reaching 1.8 updates per second for the U-Net, 1.4 for PixArt-$\alpha$ and 0.53 for SD~3.5. We train each adapter on a single RTX A4000 (16~GB), at about 1.0 updates per second.


\textbf{Decoder.} Every backbone is trained on latents from its frozen encoder, so the latent space is the contract between encoder, world model and decoder: changing the encoder would move every training target and invalidate the trained backbones, the adapters and the reconstruction upper bound.
The decoder only renders: it sits after the model and in $A$ its errors are shared by the prediction, the rendered ground truth and persistence, so fine-tuning it changes what we can see, not what the model knows.
We fine-tune the SD~1 decoder $D$ on training-map frames paired with their frozen-encoder latents, using MSE plus 0.1~LPIPS for 3,486 updates and gating it on held-out training-map frames (\appref{app:protocol}); it never sees an unseen arena.
It raises the scene reconstruction from 26.9 to 28.4~dB and lowers its LPIPS from 0.092 to 0.060, yet moves $A$ by only 0.2~dB zero-shot and 0.7~dB after adaptation, so the shift lives in the model, not the renderer; the LPIPS term is there because MSE alone more than triples scene LPIPS (0.092 to 0.339).
We use it for every pixel score of the U-Net and PixArt-$\alpha$ unless stated otherwise; SD~3.5 uses its own stock decoder.

\textbf{In-distribution performance.} On the training maps all three backbones beat persistence one tic ahead in PSNR and LPIPS and turn the right way ($A$ from \prov{+3.89} to +5.06~dB, directional 0.84 to 0.86 against 0.885 for the ground-truth frames; Table~\ref{tab:unseen}).
Four tics ahead they reach 21.33, 21.29 and \prov{21.59}~dB against persistence's 19.21~dB (stock decoders, full frame; SD~3.5 at 140k).

\textbf{Post-training.} To move a trained model to a new arena we keep its weights frozen and train a small adapter on a few episodes of that arena, then score the adapted model on held-out episodes at each update budget.
For each unseen arena we adapt the U-Net 200k EMA weights with rank-16 LoRA~\citep{hu2022lora} on every attention projection of the frozen backbone~\citep{gao2024vista} and train the control MLP, input projection and noise-bucket embedding in full (4.2M parameters, 0.49 percent)~\citep{xie2023difffit}.
We train on 8 adaptation episodes (1.1 A4000 GPU-hours per arena for 4k updates) and score the non-EMA weights at 0, 250, 500, 1k, 2k and 4k updates; forgetting ($A$ on the training maps), the directional check and a full fine-tune comparator at $2\times10^{-5}$ are in progress (\tbd{}).
An arena's budget is the first grid point at which $A$ reaches $\tfrac12(A_0+A_\text{train})$, half of its gap from its zero-shot $A_0$; arenas that do not by the final budget of the grid are right-censored~\citep{taylor2009transfer}.
Half of the gap in dB is not half of the error: the median gain of 1.8~dB removes 34 percent of the model's starting squared error (geometric mean over windows).

\begin{figure}[t]
\centering
\figslot{figures/fig1_method.pdf}{\linewidth}{1.8in}{Figure 2: method overview}
\caption{What we built: (a) the benchmark, (b) three backbones under one recipe, (c) post-training, (d) evaluation. In $A$, prediction, ground truth and persistence are all rendered by the same decoder $D$; $M$ and $G$ compare with the raw frame.}
\label{fig:method}
\end{figure}

\begin{table}[t]
\centering\footnotesize\setlength{\tabcolsep}{3.2pt}
\caption{Training maps (pooled means over windows) and unseen arenas (medians over arenas; in parentheses, arenas with $A>0$ or $M<0$), scene rows, one tic. U-Net and PixArt-$\alpha$ through the fine-tuned SD~1 decoder, SD~3.5 through its own (provisional 170k\ifdraftmarks{}, blue\fi). $A$ against the rendered ground-truth frame, the rest against the raw frame; directional pooled (ground truth 0.885, 0.892).}
\label{tab:unseen}
\begin{tabular}{llcccccc}
\toprule
Model & Maps & PSNR & LPIPS & $A$ (dB) & $M$ & $G$ (dB) & Directional \\
\midrule
Persistence & training & 20.95 & 0.239 & 0 & 0 & -- & 0 \\
 & unseen & 19.75 & 0.222 & 0 & 0 & -- & 0 \\
SD 1.4 U-Net & training & 25.20 & 0.158 & +5.06 & $-$0.081 & 3.36 & 0.855 \\
 & unseen & 22.30 & 0.303 & +2.03 (13/13) & +0.081 (0/13) & 4.96 & 0.804 \\
PixArt-$\alpha$ & training & 25.18 & 0.159 & +5.05 & $-$0.081 & 3.37 & 0.842 \\
 & unseen & 22.38 & 0.285 & +2.19 (13/13) & +0.062 (0/13) & 4.71 & 0.808 \\
SD 3.5 Medium & training & \prov{24.50} & \prov{0.146} & \prov{+3.89} & \prov{$-$0.094} & \prov{5.48} & \prov{0.850} \\
 & unseen & \prov{21.55} & \prov{0.280} & \prov{+1.50 (13/13)} & \prov{+0.059 (0/13)} & \prov{6.98} & \prov{0.793} \\
\bottomrule
\end{tabular}
\end{table}

%===============================================================================

\section{Results}
\label{sec:results}

We report what a model loses zero-shot, what post-training recovers and with what budget and what predicts how far an arena gets.

\textbf{Zero-shot: where the deficit sits.} What does a model keep and what does it lose when it meets an arena it never saw? We answer with the four scores of Section~\ref{sec:bench}. On unseen arenas the turn response survives (U-Net reversal rates 0.73 to 0.86 per arena against 0.79 to 0.91 in distribution; Table~\ref{tab:unseen}).
$A$ stays positive on every arena and backbone but falls from 5.06 to a median 2.03~dB, while the decoder's scene reconstruction PSNR is 25.2 to 29.7~dB and $G$ grows from 3.4 to a median 5.0~dB.
Nor is the perceptual sign reversal the decoder's floor: its own reconstruction LPIPS (0.05 to 0.09) is far below persistence's (0.18 to 0.28), yet the prediction's LPIPS exceeds it on 13 of 13 arenas ($M$ +0.03 to +0.17, median +0.08, against $-$0.08 in distribution; Figure~\ref{fig:shift}b).
The three backbones lose the same things (Table~\ref{tab:unseen}, Figure~\ref{fig:shift}a): the shift looks like a property of the task, not of the architecture or the scale.
What the model keeps is how actions act on the world: zero-shot on arena 7 the weapon fires and the view advances as in the ground truth (Figure~\ref{fig:teaser}) and the turn response holds; the loss sits mostly in the scene, whose rows fall from 27.5 to a median 24.0~dB on unseen arenas while the map-independent HUD rows move from 35.5 to 33.8~dB (U-Net, PSNR against the rendered ground truth, as in $A$).

\begin{figure}[t]
\centering
\figslot{figures/fig2d_family_step.pdf}{2.25in}{1.5in}{Figure 3a: family step}\hfill
\figslot{figures/tuned/fig3b_zero_shot_absolute.pdf}{3.15in}{1.5in}{Figure 3b: zero-shot absolute}
\caption{(a) The shift is a step, not a slope: every unseen arena has a lower zero-shot skill score $S_0$ than every training map for all three backbones and among the arenas $S_0$ does not follow $d$ (Spearman $-0.05$, $-0.15$, $-0.17$). (b) Every model beats persistence (grey bars) in PSNR on all 13 arenas and loses to it in LPIPS on 13 of 13. Scene rows, one tic, 95\% episode-bootstrap intervals; filled: fine-tuned SD~1 decoder (U-Net, PixArt-$\alpha$); open: SD~3.5's own (provisional 170k).}
\label{fig:shift}
\end{figure}

\textbf{Post-training: what it takes.} We next adapt a small subset of the U-Net's parameters on each unseen arena and ask what it takes to recover what was lost.
Eight episodes and 4k updates raise the median $A$ from 2.04 to 3.80~dB and on average 69 percent of each arena's gain comes within 250 updates (Figure~\ref{fig:adapt}, Table~\ref{tab:cost}).
9 of 13 arenas close half of their gap, with a median budget of 1k updates [250, ${>}$4k] and 5 arenas within 500; one reaches $A_\text{train}$.
Data matters less than the first updates: on the four arenas of the data ladder, one episode already gives most of the gain of sixteen (\appref{app:perarena-adapt}).
At 4k the prediction's LPIPS is below persistence's on 11 of 13 arenas, the other two within 0.01 and the median $G$ falls from 5.0 to the in-distribution 3.4~dB.

\textbf{What predicts the outcome.} The frozen model's zero-shot skill score $S_0$ tracks where an arena ends up, not how much it gains (exploratory Spearman $+0.90$ with $A$ at 4k, $+0.18$ with the gain, $-0.61$ with the budget); a frame distance does not. Before scoring any map we froze $d$, the sliced Wasserstein-2 distance between the motion-weighted cloud of a map's per-frame SD~1 latents and that of its nearest training map (\appref{app:distance}).
It places the training maps' held-out episodes inside the train-versus-train floor (0.03 to 0.07 against 0.02 to 0.09) and every unseen arena outside it (0.12 to 0.27), but within the 13 arenas it orders nothing (Spearman $-0.22$ with $A$, $-0.35$ with the budget); two registered transition-level distances fail even the separation (\appref{app:distance}).

\begin{figure}[t]
\centering
\figslot{figures/tuned/fig4_dots.pdf}{\linewidth}{1.4in}{Figure 4: adaptation}
\caption{Eight episodes and 4k updates lift the median $A$ from 2.04 to 3.80~dB, most of it by 250 updates (a); 9 of 13 arenas close half of their gap to $A_\text{train}$ (b; budgets above dots; grey: censored).}
\label{fig:adapt}
\end{figure}

\begin{table}[tb]
\centering\footnotesize\setlength{\tabcolsep}{2pt}
\caption{What closing the gap takes (non-EMA weights, fine-tuned decoder; brackets: 95\% arena-bootstrap intervals; ${>}$4k: censored; arena 7 is the highest-distance arena).}
\label{tab:cost}
\begin{tabular}{lccc}
\toprule
 & LoRA, 13 arenas, median & LoRA, arena 7 & Full fine-tune \\
\midrule
Arenas past half gap / $A_\text{train}$ by 4k & 9 [6, 12] / 1 of 13 & -- & \tbd{} \\
Budget to half gap / $A_\text{train}$ (updates) & 1k [250, ${>}$4k] / ${>}$4k & 250 / ${>}$4k & \tbd{} / \tbd{} \\
$A$ (dB) / $M$ / $G$ (dB) at 4k & 3.80 [3.45, 4.39] / $-$0.021 / 3.44 & 3.86 / $-$0.026 / 4.65 & \tbd{} \\
Forgetting (dB) / directional at 4k & \tbd{} / \tbd{} & \tbd{} / \tbd{} & \tbd{} \\
\bottomrule
\end{tabular}
\end{table}

%===============================================================================

\section{Future work}
\label{sec:limits}

On unseen arenas a Doom world model keeps its turn response and its action effects and loses the arena's appearance; eight episodes and a few hundred adapter updates recover more than half of its deficit on most arenas and restore a perceptual advantage over persistence on 11 of 13.
The claims cover 13 arenas of one level pack, one agent, one-tic scoring and mostly one seed and persistence is a weak reference where the camera moves fast.
One-tic quality does not guarantee stable rollouts: SD~3.5, best at one tic on the full frame with stock decoders, collapses in 23 non-EMA and 9 EMA rollouts of 256 and the EMA removes checkpoint-specific collapses, not the failure mode (\appref{app:closedloop}).
Next are other level packs and other games, to see whether the same recipe transfers; detecting a shift from the model's own zero-shot skill score so that it repairs itself; closed-loop evaluation of the adapted models; and agents trained inside the adapted model.
The pattern we measure, control kept, appearance lost, most of it relearned from a few episodes, is what a deployed world model would need to detect and repair on its own and DoomShift is a testbed for that loop.
\label{body-end}

%===============================================================================

\clearpage
\acknowledgments{\tbd}

\bibliography{refs}

\ifwithappendix
\clearpage
\appendix
% appendix.tex: input by main.tex after the bibliography when \withappendixtrue.
% The PhysWM call limits the paper to "up to 4 pages (excluding references)" and says nothing about an
% appendix; the OpenReview form takes one PDF. Whether this rides along in the submission PDF is Rohan's call.
% Marks as in main.tex: \tbd pending, \todo{} pending text, \prov{} provisional (replaced by the fresh-set
% rescore, the tuned decoder or the final checkpoints); each paragraph that carries one names its source.

\section{Protocol details}
\label{app:protocol}

% sources: docs/RESEARCHER_DOSSIER.md 2.1 to 2.3 (recorder, executed controls), release/DENSE_CORPUS.md (map split); design page section 3 (fresh set); opus-distance-view section 1 (the arenas 6 to 8 batch)
% depends: fresh 24-episode recording and SD 1 encode of the 13 unseen arenas (Sep 26 to 27)
\textbf{Recording.} \texttt{record\_arnold.py} replaces Arnold's evaluation loop with one that steps the engine one tic at a time and stores each frame, then the control it is about to send, then steps. The raw \texttt{buttons} field is the list Arnold \emph{requested}; the engine executes its first 19 entries. Arnold's action builder appends ten weapon-selection names to that list at every game start, so a weapon request lands inside the executed range only in a recorder worker's first episode, and later weapon requests are silently dropped while movement and firing still execute. Requested and executed controls agree on 89 percent of tics, almost all of the rest are dropped weapon requests, and 0.15 percent are anti-stuck overrides. The sidecars store the executed 19-bit string, which is what the models condition on. For the fresh unseen-arena set, \prov{328} episodes were recorded by 16 workers under one pinned protocol (same Arnold build, 8 bots, 150~s, every tic, seeds declared in the corpus README and disjoint from every earlier batch), and each worker's first episode was set aside, leaving \prov{24} per arena. The earlier batches, which the pre-fresh-set numbers in this appendix use, held 4 episodes on arenas 1 and 9 to 15, 10 on 16 and 17 and 20 on 6 to 8; arenas 6 to 8 came from a separate batch that scored lower in decoded advantage (\prov{$p=0.03$}), a batch-map confound the single-protocol set removes.

\textbf{Windows.} A one-tic window is 32 consecutive recorded tics plus the target, inside one life on one map; deaths cut windows, which keeps 85 (arena 1) to 99 percent of candidate starts. The in-domain reads use 512 windows drawn with seed 0 from validation episodes 6000 to 6099 (25 per training map); the per-arena reads use 256 windows per map.

\begin{table}[ht]
% sources: docs/RESEARCHER_DOSSIER.md section 3 table (parameters, memory, throughput), backbones.py (control paths)
\centering\footnotesize\setlength{\tabcolsep}{3pt}
\caption{The three rows under one recipe. Memory and throughput at batch 32 on one A6000. Control tokens enter through cross-attention (U-Net, width 768; PixArt, 4,096) or joint attention (SD~3.5, 4,096, plus a pooled slot carrying the newest control).}
\label{tab:recipe}
\begin{tabular}{lccc}
\toprule
 & SD 1.4 U-Net & PixArt-$\alpha$ 512 & SD 3.5 Medium \\
\midrule
Architecture & U-Net & DiT, cross-attention & MMDiT, joint attention \\
Latent space & SD 1.x, $4\times32\times40$ & SD 1.x, $4\times32\times40$ & SD 3.5, $16\times32\times40$ \\
Parameters & 860.5M & 628M & 2,271.7M \\
Memory; checkpointing & 23.3 GB; off & 26.8 GB; off & 36.8 GB; on \\
Updates per second & 1.78 to 1.83 & about 1.4 & 0.50 to 0.55 \\
\bottomrule
\end{tabular}
\end{table}

\textbf{Recipe notes.} The 1,000-step linear schedule ($\beta$ from $10^{-4}$ to 0.02) is neither backbone's native schedule: 273 of its timesteps are noisier than anything SD~1.4 saw in pretraining, and SD~3.5 was pretrained with rectified flow. Context noise augmentation blends $\sqrt{1-q}\,c+\sqrt{q}\,\epsilon$ with one level $q$ per example drawn below 0.7 and embedded as one of 10 buckets; inference uses clean context and bucket 0. Fused AdamW runs at $5\times10^{-5}$ after 2,000 linear warmup updates, with no weight decay, gradient clipping at 1.0, bf16 autocast over fp32 parameters, seed 0 and no action dropout. The EMA decays at 0.9999 per update, applied as $0.9999^8$ every eighth update, in fp32. Training sees windows from episodes 0 to 1,999 (500 per training map, 9.66M valid windows); 200k updates at batch 32 cover about two thirds of one pass. The decoder tune (\texttt{finetune\_decoder.py}) trains the SD~1.x and SD~3.5 decoders with MSE on training-map frames with the encoder frozen, validates on episodes 6000 to 6099, and keeps the terminal checkpoint; the stock decoders saw no Doom frame.

% sources: dossier 4.2 (U-Net 200k, live weights, h1, 512 training-map validation windows)
\textbf{Sampler.} Fewer DDIM steps land nearer the conditional mean, which wins PSNR and loses LPIPS. U-Net 200k, live weights, one tic, training-map windows, linear spacing: 4 steps 22.35~dB / 0.272, 8 steps 22.34 / 0.199, 10 steps 22.29 / 0.186, 16 steps 22.20 / 0.171, 50 steps 22.04 / 0.162. Every number in the paper uses 10 steps.

\section{Per-arena zero-shot scores}
\label{app:perarena}

\begin{table}[ht]
% depends: fresh-set rescore of the 13 arenas and the training maps (tuned SD 1 decoder, scene rows, per-window CIs by bootstrap over episodes); D recomputed on the fresh set; U-Net 30-map directional (15 of 30 maps done; arena 6 reads 0.789 against ref 0.867)
% rows generated from results/distance_study/scores/040-unet-nexttic_snap_0200000_ema_ddim10/*_h{1,4}/metrics.json and distances_sd1.json (primary entries), not typed
\centering\footnotesize\setlength{\tabcolsep}{3.5pt}
\caption{Per arena, U-Net 200k EMA, 256 windows per map, stock decoder, full frame, sorted by $D$. Persistence and raw gain are against the raw last frame; $A$, $M$ and $G$ as in Table~\ref{tab:unseen}. Blue: pre-fresh-set episodes. $^\ast$pooled over the four training maps.}
\label{tab:perarena}
\begin{tabular}{rlccccccccc}
\toprule
Arena & Role & $D$ & Persist. & Ceiling & $A$, 1 tic & $A$, 4 tics & $M$ & $G$ & Raw gain & Directional \\
\midrule
5 & training & \prov{0.026} & \prov{19.76} & \prov{23.38} & \prov{+3.70} & \prov{+3.22} & \prov{$-$0.029} & \prov{1.84} & \prov{+1.79} & 0.867$^\ast$ \\
3 & training & \prov{0.034} & \prov{21.37} & \prov{23.71} & \prov{+3.21} & \prov{+4.01} & \prov{$-$0.006} & \prov{1.62} & \prov{+0.72} & 0.867$^\ast$ \\
2 & training & \prov{0.065} & \prov{21.92} & \prov{23.80} & \prov{+3.23} & \prov{+3.23} & \prov{$-$0.016} & \prov{1.37} & \prov{+0.51} & 0.867$^\ast$ \\
4 & training & \prov{0.069} & \prov{22.70} & \prov{24.74} & \prov{+4.26} & \prov{+4.98} & \prov{$-$0.056} & \prov{1.36} & \prov{+0.68} & 0.867$^\ast$ \\
\midrule
8 & unseen & \prov{0.115} & \prov{23.20} & \prov{24.33} & \prov{+1.29} & \prov{+1.05} & \prov{+0.077} & \prov{2.05} & \prov{$-$0.92} & \tbd \\
15 & unseen & \prov{0.127} & \prov{19.85} & \prov{23.75} & \prov{+2.68} & \prov{+1.14} & \prov{+0.079} & \prov{2.61} & \prov{+1.30} & \tbd \\
12 & unseen & \prov{0.128} & \prov{21.69} & \prov{23.39} & \prov{+1.07} & \prov{+0.96} & \prov{+0.104} & \prov{2.01} & \prov{$-$0.31} & \tbd \\
17 & unseen & \prov{0.130} & \prov{19.37} & \prov{22.84} & \prov{+1.30} & \prov{+0.79} & \prov{+0.155} & \prov{2.94} & \prov{+0.53} & \tbd \\
14 & unseen & \prov{0.146} & \prov{21.75} & \prov{23.87} & \prov{+1.33} & \prov{+0.35} & \prov{+0.109} & \prov{2.39} & \prov{$-$0.28} & \tbd \\
11 & unseen & \prov{0.150} & \prov{20.83} & \prov{24.09} & \prov{+1.80} & \prov{+0.95} & \prov{+0.080} & \prov{2.91} & \prov{+0.35} & \tbd \\
13 & unseen & \prov{0.156} & \prov{19.37} & \prov{22.69} & \prov{+1.80} & \prov{+0.90} & \prov{+0.125} & \prov{2.53} & \prov{+0.79} & \tbd \\
10 & unseen & \prov{0.162} & \prov{22.60} & \prov{24.20} & \prov{+1.67} & \prov{+1.93} & \prov{+0.069} & \prov{1.96} & \prov{$-$0.36} & \tbd \\
16 & unseen & \prov{0.164} & \prov{20.03} & \prov{23.03} & \prov{+1.92} & \prov{+1.11} & \prov{+0.163} & \prov{2.23} & \prov{+0.77} & \tbd \\
1 & unseen & \prov{0.179} & \prov{21.93} & \prov{23.74} & \prov{+1.09} & \prov{+0.49} & \prov{+0.138} & \prov{2.04} & \prov{$-$0.23} & \tbd \\
9 & unseen & \prov{0.180} & \prov{22.58} & \prov{24.63} & \prov{+2.05} & \prov{+1.64} & \prov{+0.063} & \prov{2.11} & \prov{$-$0.06} & \tbd \\
6 & unseen & \prov{0.188} & \prov{20.72} & \prov{22.85} & \prov{+0.72} & \prov{$-$0.06} & \prov{+0.122} & \prov{2.44} & \prov{$-$0.31} & \prov{0.79} \\
7 & unseen & \prov{0.282} & \prov{19.53} & \prov{23.58} & \prov{+0.88} & \prov{$-$0.76} & \prov{+0.084} & \prov{4.00} & \prov{+0.05} & \tbd \\
\bottomrule
\end{tabular}
\end{table}

% recomputed from the same files (Spearman over the 13 unseen arenas); HUD share: (32/240) x HUD-row MSE over full-frame MSE of the ceiling, from hud_vae_psnr and vae_psnr
% depends: fresh-set rescore (every correlation below is recomputed on it)
\textbf{What the variation between arenas follows.} Across the 13 unseen arenas, raw gain over persistence tracks persistence PSNR (Spearman \prov{$-0.77$}) and the model's absolute PSNR tracks it at \prov{0.95}, while $A$ does not (\prov{$-0.05$}): the model's pixel error has a floor set by the decoder and by blur, and raw persistence has none, so on static footage raw persistence wins however good the latent prediction is. The gap $G$ tracks persistence too (\prov{$-0.80$}), which is why it is a table column and not a cost axis. The 32 HUD rows, which the stock decoder renders at about 18~dB and persistence copies exactly, carry 39 to 68 percent of the ceiling's squared error on the 17 arenas.

\section{The distance}
\label{app:distance}

% sources: distance-study-design-2026-09-24.md sections 2 and 5; dossier 7.2 to 7.5; results/distance_study/distances_sd1.json (config, floor, checks) and figure_unet_h1_amended/stats.json (primary, within_cluster, amended_assessment)
% depends: D recomputed on the fresh set with the frozen code; the transition-level distance frozen before any adapted score
\textbf{Definition, frozen Sep 25 before any map was scored.} A map is a cloud of 1,000 frames: 4 episodes $\times$ 250 frames at least 32 tics into a life, averaged over ten random 4-episode subsets where a map has more episodes. Features are SD~1.x latents with the padding rows removed, pooled $5\times5$ to 192 dimensions. Each frame is weighted by its latent motion $\lVert z_t-z_{t-1}\rVert$ plus a floor that gives the quietest half of the cloud 25 percent of the weight. The reference is 50 seeded training episodes per training map, 250 frames each, 25 per motion decile. $D(m)=\min_k \mathrm{SW}_2(\text{cloud}_m,\text{reference}_k)$ over the four training maps, with 1,000 projection directions shared by every distance; the nearest map, not the pooled corpus, because transport to a mixture charges a seen map for not resembling the other three. A disjoint second reference draw gives the train-versus-train floor, 0.022 to 0.093. Pixel-space and SD~3.5-space versions rank the 30 maps of the pre-registered set like the primary (Kendall $\tau$ 0.70 and 0.76).

\begin{table}[ht]
% sources: dossier 7.4 and results/distance_study/distances_sd1.json checks; the amended reading is figure_unet_h1_amended/stats.json amended_assessment
\centering\footnotesize\setlength{\tabcolsep}{4pt}
\caption{The pre-declared validation checks of $D$ on the 30-map set (SD~1.x space). The pre-declared gate required all five; it failed, and the pre-declared verdict ``does not support'' stands. The dated amendment separates checks of the measurement (i, iii, v) from an assumption about the maps (iv) and a subset tolerance stricter than the estimator reported (ii).}
\label{tab:checks}
\begin{tabular}{clc}
\toprule
Check & What it tests & Result \\
\midrule
(i) & the four training maps sit inside the floor band & pass (0.026 to 0.069) \\
(ii) & disjoint 4-episode subsets agree within 10\% & fail (8 campaign maps, 0.108 to 0.288) \\
(iii) & the two reference draws rank maps alike ($\rho\ge0.95$) & pass (0.996) \\
(iv) & arenas 6 to 8 sit nearer than every campaign map & fail (arena 7 at 0.282) \\
(v) & $|\rho(D,\text{Kish }n_\text{eff})|<0.4$ & pass ($-$0.022) \\
\bottomrule
\end{tabular}
\end{table}

% sources: results/distance_study/figure_unet_h1/stats.json (primary, leave-one-out, within_cluster, secondary.gain_h4) and figure_unet_h1_amended/stats.json (structural_controls, summary)
% depends: the Spiderman rerun of the structural controls with the per-window files (the local amended CIs resampled without them); the within-unseen partial is recomputed on the fresh set
\textbf{The pre-registered test.} On 30 maps (the 17 arenas and 13 campaign maps of another WAD that the paper otherwise leaves in the dataset release), the partial Spearman of the U-Net's one-tic raw gain on $D$, controlling for persistence PSNR, is $-0.73$ [$-0.84$, $-0.32$] (permutation $p=10^{-4}$), negative in all 30 leave-one-map-out refits and $-0.80$ at four tics. Within the 17 arenas it is $-0.57$, within the campaign maps $-0.38$ with an interval that crosses zero, and with a training-map indicator and a family indicator as covariates $-0.42$ [\prov{$-0.63$}, \prov{$+0.13$}]. Within the 13 unseen arenas it is \prov{$-0.12$}. The dated amended assessment reads ``qualified evidence for the pooled U-Net association; original gate failed; unseen-only trend inconclusive; cross-backbone replication pending''.

% depends: the transition-distance design round (Astra's directed coverage of transition windows; Opus's kNN transfer gap), validated on zero-shot data and frozen before any 13-arena adapted score; computed on the fresh set
\textbf{The transition-level distance.} $D$ reads single frames; the model predicts a transition from 32 frames and a control. The second axis of the adaptation study is a directed distance on transition windows (state, recent motion, next change and the executed controls), from the target's windows to the nearest training windows, chosen between two candidates and frozen before any adapted score: \todo{the chosen definition and its validation: floor on disjoint source episodes, a rise under shuffled transitions and shuffled controls, split-half map reliability}. Directed coverage predicted transfer better than symmetric transport in two studies~\citep{mensink2021factors,westny2026latent}.

\section{Closed loop}
\label{app:closedloop}

% counts recomputed from results/sd35_stability/collapse_rates_50k_130k.md; mechanism from dossier 6.2; the live-minus-EMA interval is Astra's cluster bootstrap over the 16 windows (astra-review-2026-09-26.md)
% depends: PixArt 200k final read (two-seed rollouts) and SD 3.5 200k read add rows
\textbf{The failure.} In SD~3.5's collapsed rollouts the latent's overall scale stays inside the real range while one channel's per-frame mean is captured at a fixed value: channel 13, whose mean sits near $+0.45$ in real footage, drops to $-1.39$ to $-1.50$ (17 of the 22 collapsed rollouts at 70k, both window sets) or $-2.60$ to $-2.73$ (the blank frames). Later EMA excursions shift channels 15 and 8 as well, so the absorbing states are a family of per-channel mean shifts. Noise augmentation trains against isotropic corruption of the context, not a structured shift in one direction of latent space, and teacher-forced reads never feed a prediction back, so neither guards against it.

\begin{table}[ht]
\centering\footnotesize\setlength{\tabcolsep}{5pt}
\caption{SD~3.5 collapse events on the same 16 seed-0 windows at every 5k from 55k to 130k (256 rollouts per weight set). (a) a frame under 10~dB; (b) under 12~dB at tic 256; (c) channel 13's mean below $-0.9$ at tic 256 (real range $-0.44$ to $0.80$). Live minus EMA for (a): 5.5 points [2.7, 7.8], bootstrapped over the 16 windows.}
\label{tab:collapse}
\begin{tabular}{lcccccc}
\toprule
 & \multicolumn{3}{c}{55k to 130k} & \multicolumn{3}{c}{without 70k} \\
Weights & (a) & (b) & (c) & (a) & (b) & (c) \\
\midrule
Live & 23 / 256 & 17 / 256 & 14 / 256 & 11 / 240 & 6 / 240 & 4 / 240 \\
EMA & 9 / 256 & 4 / 256 & 4 / 256 & 8 / 240 & 3 / 240 & 4 / 240 \\
\bottomrule
\end{tabular}
\end{table}

\begin{figure}[ht]
% exists: paper/figures/sd35_70k_live_vs_ema_rollout_strip.jpg (tools/collapse_strip.py at 42a4688)
\centering
\figslot{figures/sd35_70k_live_vs_ema_rollout_strip.jpg}{\linewidth}{3.2in}{SD 3.5 70k live versus EMA rollout strip}
\caption{SD~3.5 at 70k updates, two validation windows (map 2 above, map 5 below); rows are the recorded frames, the live-weight rollout and the EMA rollout at tics 1 to 256. The live rollouts flood to flat blue (channel 13 near $-2.6$) or grey (near $-1.4$) while the HUD, weapon and crosshair persist; the EMA rollouts drift from the true trajectory but stay coherent scenes.}
\label{fig:collapse}
\end{figure}

% sources: dossier 5.1 (U-Net 200k rollouts), RESEARCH_CONTEXT 2026-09-26 12:10 (PixArt 150k) and 17:10 (SD 3.5 140k)
% depends: PixArt 200k and SD 3.5 200k rollout reads
\textbf{EMA against live, all rows.} At 256 tics on the seed-0 windows the EMA and live weights read 17.81 and 15.87~dB for the U-Net at 200k, \prov{17.41} and \prov{15.97} for PixArt at 150k and \prov{18.03} and \prov{17.53} for SD~3.5 at 140k, against 17.52 for copy-seed. The EMA led the live weights at every recorded teacher-forced read once its start weights had washed out (by 40k for the U-Net). The adaptation curves score the live adapter weights, because an EMA at 0.999 still carries the step-0 model at 0.78 after 250 updates.

\section{Latent skill, and LoRA against full fine-tuning}
\label{app:skill}

% depends: latent-ratio rescore of the 13 arenas and the training maps (eval_tf.py latent_mse_ratio per window, fresh set); 13-arena LoRA runs; full fine-tune comparator on arena 7
The latent skill $S=10\,\overline{\log_{10}(\text{copy-last latent MSE}/\text{model latent MSE})}$ is a geometric mean over windows, so near-static windows cannot dominate it, and no decoder touches it. It compares predictions only within one backbone, because the latent spaces differ. Zero-shot $S$ is \tbd{}~dB on the training maps and \tbd{}~dB on the unseen arenas. Because a full fine-tune also moves the input projection and every layer the decoder's input depends on, the LoRA-against-full comparison is read in $S$ as well: on arena 7, $S$ at 4,000 updates is \tbd{} for LoRA and \tbd{} for the full fine-tune.

\section{Per-arena adaptation}
\label{app:perarena-adapt}

\begin{table}[ht]
% depends: 13-arena U-Net LoRA runs (score_adapt.py score and cost; live weights), forgetting and directional guards; D and step-0 A are today's provisional values and are replaced by the fresh-set rescore
\centering\footnotesize\setlength{\tabcolsep}{3.5pt}
\caption{Per unseen arena, sorted by $D$: step-0 decoded advantage, the first grid budget at which $A$ reaches 25, 50 and 100 percent of the training maps' advantage (${>}$4k: censored), $A$ and $M$ at 4,000 updates, forgetting (change in the training maps' $A$) and the directional score at 4,000.}
\label{tab:perarena-adapt}
\begin{tabular}{rccccccccc}
\toprule
Arena & $D$ & $A$ at 0 & 25\% & 50\% & 100\% & $A$ at 4k & $M$ at 4k & Forgetting & Directional \\
\midrule
8 & \prov{0.115} & \prov{+1.29} & \tbd & \tbd & \tbd & \tbd & \tbd & \tbd & \tbd \\
15 & \prov{0.127} & \prov{+2.68} & \tbd & \tbd & \tbd & \tbd & \tbd & \tbd & \tbd \\
12 & \prov{0.128} & \prov{+1.07} & \tbd & \tbd & \tbd & \tbd & \tbd & \tbd & \tbd \\
17 & \prov{0.130} & \prov{+1.30} & \tbd & \tbd & \tbd & \tbd & \tbd & \tbd & \tbd \\
14 & \prov{0.146} & \prov{+1.33} & \tbd & \tbd & \tbd & \tbd & \tbd & \tbd & \tbd \\
11 & \prov{0.150} & \prov{+1.80} & \tbd & \tbd & \tbd & \tbd & \tbd & \tbd & \tbd \\
13 & \prov{0.156} & \prov{+1.80} & \tbd & \tbd & \tbd & \tbd & \tbd & \tbd & \tbd \\
10 & \prov{0.162} & \prov{+1.67} & \tbd & \tbd & \tbd & \tbd & \tbd & \tbd & \tbd \\
16 & \prov{0.164} & \prov{+1.92} & \tbd & \tbd & \tbd & \tbd & \tbd & \tbd & \tbd \\
1 & \prov{0.179} & \prov{+1.09} & \tbd & \tbd & \tbd & \tbd & \tbd & \tbd & \tbd \\
9 & \prov{0.180} & \prov{+2.05} & \tbd & \tbd & \tbd & \tbd & \tbd & \tbd & \tbd \\
6 & \prov{0.188} & \prov{+0.72} & \tbd & \tbd & \tbd & \tbd & \tbd & \tbd & \tbd \\
7 & \prov{0.282} & \prov{+0.88} & \tbd & \tbd & \tbd & \tbd & \tbd & \tbd & \tbd \\
\bottomrule
\end{tabular}
\end{table}

\begin{figure}[ht]
% depends: 13-arena U-Net LoRA runs and score_adapt.py rows (live weights, 256 held-out windows per arena); the line from step-0 training-map predictions through the tuned SD 1 decoder; the full fine-tune comparator on arena 7
\centering
\figslot{figures/fig5_adaptation_margin.pdf}{0.6\linewidth}{1.8in}{Figure 5: $M$ against LoRA updates, one line per arena}
\caption{The second panel of the adaptation study: perceptual margin $M$ against LoRA updates, one line per unseen arena coloured by $D$, with zero (ties copy-last) and the training maps' margin; open markers: arenas censored at 4,000 updates. Black: full fine-tune, arena 7.}
\label{fig:adapt-margin}
\end{figure}

\section{Public Doom world models}
\label{app:public}

% sources: main-, astra- and opus-decoder-baselines-view-2026-09-26.md (Hugging Face API file lists and READMEs checked Sep 26)
Two open GameNGen reproductions (Stiegler's and Taketani's) release weights; both train on one map (\texttt{deathmatch\_simple}) at 320$\times$240 every tic, with 18 and 12 discrete actions that do not cover Arnold's combined controls. DIAMOND's released model is CS:GO, Oasis is a gated Minecraft model, and PlayGen and MultiGen release no Doom weights. Scoring a one-map public model on our arenas would measure its own transfer to Arnold's footage, confounded by map, assets, behaviour and recipe, not compare recipes; we leave it to a later version, scored home and away with only the windows whose control history maps exactly.

\fi

\end{document}
